\chapter*{学生实验}\markboth{学生实验}{学生实验}
\stepcounter{chapter}
\ctexset{section/name={实验,}}
\addcontentsline{toc}{chapter}{学生实验}
\renewcommand{\thefigure}{\arabic{figure}}
\renewcommand{\thetable}{\arabic{table}}

\section{氨的制备和性质\texorpdfstring{\quad}{ }铵离子的检验}
\subsection*{实验目的}
\begin{enumerate}[1.]
  \item 学会实验室制取氨的方法；
  \item 认识氨的物理性质和化学性质；
  \item 学会检验氨和铵离子的方法。
\end{enumerate}
\subsection*{实验用品}
试管、铁架台、酒精灯、玻璃棒、药匙、玻璃片、研钵、带导管的塞子、胶头滴管、水槽、棉花。

氯化铵、硫酸铵、硝酸铵、氢氧化钙、氢氧化钠溶液、浓盐酸、浓硫酸、浓硝酸、红色石蕊试纸、酚酞试液。
\subsection*{实验步骤}
\begin{enumerate}[1.]
  \item 氨的制取
  \begin{enumerate}[(1)]
    \item 取氯化铵和氢氧化钙各一药匙，放在研钵里，用玻璃棒充分拌和，有什么气味发生？说明发生了什么反应？写出化学方程式。
    \item 把上面的混和物放到试管里，用带导管的塞子塞住试管口，并固定在铁架台上。导管的另一头向上伸入另一个干燥的倒置的试管里（见\cref{fig:a1}）。
    \begin{figure}
      \includegraphics{A-1.pdf}
      \caption{氨的制取}\label{fig:a1}
    \end{figure}
    \item 用小火加热试管，便有氨气生成。
  \end{enumerate}
  \item 氨的性质
  \begin{enumerate}[(1)]
    \item 在收集氨气的试管口处，放一条润湿的红色石蕊试纸，待红色石蕊试纸变蓝时，即说明试管中的氨气已经收集满了，这时停止加热。把倒立的试管轻轻拿下，并用拇指堵住管口。注意观察氨气的颜色、气味、状态。
    \item 将上述充满氨气的试管，管口向下倒拿着放到水槽的水里去（\cref{fig:a2}）。把拇指放开。看到什么现象？为什么？
    \begin{figure}
      \begin{minipage}[b]{0.48\linewidth}\centering
        \includegraphics{A-2.pdf}
        \caption{氨在水中的溶解}\label{fig:a2}
      \end{minipage}
      \begin{minipage}[b]{0.48\linewidth}\centering
        \includegraphics{A-3.pdf}
        \caption{氨与酸的反应}\label{fig:a3}
      \end{minipage}
    \end{figure}
    \item 当水进入试管后，在水面下用拇指堵住试管，把试管从水里取出，管口向上，振荡试管，并向溶液里滴入几滴酚酞试液，观察有什么现象？
    \item 将实验步骤 1 里制取氨的试管按\cref{fig:a3} 装好。用玻璃棒蘸浓硫酸、浓硝酸和浓盐酸，在玻璃片上的不同地方各滴一滴（玻璃棒每取过一种酸后，要用水洗干净后，再用来取另一种酸）。然后加热氯化铵和氢氧化钙的混和物，当有氨放出时，移动玻璃片，使导管口依次对着三滴不同的酸。这时各有什么现象产生？为什么有的酸滴上冒白烟？玻璃片上生成的三种白色物质各是什么？写出三个反应的化学方程式。
  \end{enumerate}
  \item 铵离子的检验
  取少量氯化铵、硝酸铵和硫酸铵晶体，分别放在三个试管里，然后用胶头滴管分别滴入少量氢氧化钠溶液，加热试管，再把湿润的红色石蕊试纸放在管口处，观察试纸的颜色有什么变化？写出化学方程式，根据这个实验可以得出什么结论？
\end{enumerate}
\begin{Theorem}{问题和讨论}
制取和收集氨气时应注意哪些问题？
\end{Theorem}



\section{硝酸和硝酸盐的性质}
\subsection*{实验目的}
认识硝酸和硝酸盐的特性。

\subsection*{实验用品}
试管、酒精灯、试管夹、胶头滴管、烧瓶、铁架台、导管、水槽、小木条、橡皮塞。

浓硝酸、稀硝酸（$1:2$）、浓硫酸、硝酸铜、硝酸钠、硝酸铵、铜片、石蕊试液、蓝色石蕊试纸。
\subsection*{实验步骤}
\begin{enumerate}[1.]
  \item 硝酸的性质
  \begin{enumerate}[(1)]
    \item 浓、稀硝酸对石蕊试液的作用\quad 取两个试管，分别注入浓、稀硝酸各 \qtyrange{1}{2}{ml}，再分别滴入几滴石蕊试液，微热，看到石蕊试液在稀硝酸中变红，而在浓硝酸中褪色。为什么？
    \item 浓硝酸跟铜的作用\quad 在一个试管中放入一片铜片，再滴入几滴浓硝酸，观察放出气体的颜色。
    
    向这个试管里加入 \qty{5}{ml} 水，观察溶液的颜色。

    上述现象说明有什么物质生成？写出浓硝酸跟铜反应的化学方程式。
    \item 稀硝酸跟铜的作用\quad  照\cref{fig:a4} 把仪器安装好，并检查这一装置的气密性。
    
    在烧瓶里放 \numrange{3}{4} 片铜片，再倒入稀硝酸到浸没铜片为止。把烧瓶加热，用排水取气法收集一试管一氧化氮。

    \begin{figure}
      \includegraphics{A-4.pdf}
      \caption{\ce{NO} 的生成}\label{fig:a4}
    \end{figure}

    注意观察一氧化氮的颜色。

    在水面下用拇指堵住试管口，把试管从水中取出，并使管口向上，放开拇指，观察气体颜色有什么变化？等气体全部变色后，再向试管中加入 \qty{2}{ml} 水，用塞子塞佐。振荡，使气体溶解在水中。用玻璃棒蘸取试管里的溶液，滴在蓝色石蕊试纸上，观察有什么变化？写出铜跟稀硝酸反应的化学方程式。
  \end{enumerate}
  \item 硝酸盐的性质
  \begin{enumerate}[(1)]
    \item 把少量预先烘干的硝酸铜粉末，放入干燥的试管里，加热，观察发生的变化和放出的气体的颜色。把带火星的木条仲进试管里，有什么现象发生？写出反应的化学方程式。
    \item 把少量固态硝酸钠、硝酸铜、硝酸铵分别放进三个试管里，各加入少量浓硫酸，再向每个试管里加入一小片铜。加热试管，观察有什么现象发生。写出有关反应的化学方程式。
  \end{enumerate}
\end{enumerate}
\begin{Theorem}{问题和讨论}
硝酸跟金属的反应和盐酸跟金属的反应的现象有什么不同？本质有什么区别？
\end{Theorem}


\section{化学反应速度\texorpdfstring{\quad}{ }化学平衡}
\subsection*{实验目的}
\begin{enumerate}[1.]
  \item 巩固浓度、温度和催化剂对反应速度的影响的知识；
  \item 巩固浓度和温度对化学平衡的影响的知识。
\end{enumerate}
\subsection*{实验用品}
试管、烧杯、胶头滴管、橡皮塞、电钟（或秒表）、酒精灯、温度计、铁架台、药匙、木条。

3\% 硫代硫酸钠溶液、硫酸（$1:5$）、3\% 过氧化氢溶液、氯化铁溶液、硫氰化铵溶液、两个装有二氧化氮的大试管\footnote{\ce{NO2} 气体，可在课前由教师组织部分学生制好，并装在试管里，用塞子塞好，留待实验课使用。}。二氧化锰。

\subsection*{实验步骤}
\begin{enumerate}[1.]
  \item 浓度、温度和催化剂对化学反应速度的影响
  \begin{enumerate}[(1)]
    \item 浓度对化学反应速度的影响
    
    取三个试管，分别编为 1、2、3 号，并按下表规定的数量分别加入硫代硫酸钠溶液和蒸馏水，摇匀后，把试管放在一张有字的纸前，这时隔着试管可以清楚地看到字迹。然后再加入硫酸。同时从加入第一滴硫酸时开始记录时间，到溶液出现的浑浊使试管后面的字迹看不见时，停止记时。将记录的时间填入下表。
    \begin{tablehere}
      \begin{tblr}{colspec={ccccX[c]},hline{2}=0.8pt}
        编号 & {加 3\% \ce{Na2S2O3}\\ （\unit{ml}）}& {加 \ce{H2O}\\（\unit{ml}）}& {加 \ce{H2SO4}（$1:5$）\\（滴）}& 出现浑浊所需时间（\unit{s}）\\
        1 &  5 & 5 & 10 & \\
        2 &  7 & 3 & 10 & \\
        3 & 10 & 0 & 10 & \\
      \end{tblr}
    \end{tablehere}
    \item 温度对化学反应速度的影响
    
    取三个试管，分别加入硫代硫酸钠溶液 \qty{5}{ml}，然后在室温条件下，向一个试管中滴入 5 滴硫酸，并记录出现浑浊所需的时间。再把另两个试管放在水浴中加热，使两个试管的温度分别为高于室温 \qty{10}{\celsius} 和高于室温 \qty{20}{\celsius}，再分别向这两个试管中滴入 5 滴硫酸，再用与实验步骤 1.(1) 相同的方法记录出现浑浊所需的时间。
    \begin{tablehere}
      \begin{tblr}{colspec={ccX[c]cc},hline{2}=0.8pt}
        编号 & 加 3\% \ce{Na2S2O3}（\unit{ml}）& 加 \ce{H2SO4}（$1:5$）（滴）& 温度（\unit{\celsius}） & 时间（\unit{s}）\\
        1 &  5 & 5 & 室温 & \\
        2 &  5 & 5 & 室温 + 10 & \\
        3 &  5 & 5 & 室温 + 20 & \\
      \end{tblr}
    \end{tablehere}
    \item 催化剂对化学反应速度的影响
    
    在一试管中加入约 \qty{2}{ml} 3\% \ce{H2O2} 溶液，观察是否有气泡发生。然后再加入少量二氧化锰粉末，再观察是否有气泡发生，并用带火星的木条检验产生的气体。
  \end{enumerate}
  \item 浓度和温度对化学平衡的影响
  \begin{enumerate}[(1)]
    \item 浓度对化学平衡的影响
    \begin{enumerate}[label=\protect\cnum{\arabic*}]
      \item 在小烧杯中，加入氯化铁溶液和硫氰化铵溶液各 \qty{10}{ml}，混和后得到红色溶液。把这溶液分别平均地倒入三个试管里。
      \item 向第一个试管里加入氯化铁溶液，向第二个试管中加入硫氰化铵溶液。分别将上述二个试管与第三个试管比较，观察溶液颜色的变化。
    \end{enumerate}

    根据溶液颜色的变化，说明浓度对化学平衡的影响。
    \item 温度对化学平衡的影响\quad 取两个带塞的大试管，里面装有达到平衡的 \ce{NO2} 和 \ce{N2O4} 气体。
    \[\schemestart 
    \chemname{\ce{2NO2}}{\footnotesize （红棕色）}
    \arrow(.mid east--.mid west){<=>}
    \chemname{\ce{N2O4}}{\footnotesize （无色）}
    \+ \qty{13.6}{kCal}
    \schemestop\chemnameinit{}\]

    将一个试管浸在盛有热水的烧杯里，另一个试管浸在盛有冷水的烧杯里，比较两试管里气体的颜色，并说明温度对化学平衡的影响。
  \end{enumerate} 
\end{enumerate}
\begin{Theorem}{问题和讨论}
小结影响化学反应速度和化学平衡的条件。
\end{Theorem}


\section{实验习题}
\subsection*{实验目的}
巩固已学知识，培养实验能力。

\subsection*{实验习题}
\begin{enumerate}[1.]
  \item 用实验方法鉴别：
  \begin{tasks}[item-indent=5em,label=\Alph*.,label-width=15pt](2)
    \task 浓硝酸和稀硝酸。
    \task 浓硫酸和稀硫酸。
    \task 盐酸和硝酸。
  \end{tasks}
  \item 有四种白色粉末，它们是硫酸铵、硝酸钠、氯化铵和氯化钠，通过实验把它们鉴别出来。
  \item 怎样通过简单的实验证明碳酸氢铵加热分解的产物是什么。
  \item 用什么简单方法除去混在氯化钠晶体里的少量氯化铵？怎样检验是不是除净了？
  \item 只用一种试剂，鉴别下列溶液。
  \[\ce{Na2CO3},\quad\ce{AgNO3},\quad\ce{NH4Cl}\]
  \item 用浓度都是 $0.02M$ 的 \ce{Ca(OH)2} 和 \ce{H3PO4} 依次制取 \ce{Ca3(PO4)2}、\ce{CaHPO4}、\ce{Ca(H2PO4)2}。如果 \ce{Ca(OH)2} 的用量为 \qty{1.5}{ml}，计算 \ce{H3PO4} 的用量。根据计算结果，在盛有 \qty{1.5}{ml} \ce{Ca(OH)2} 溶液的试管里，先加入一定量的 \ce{H3PO4}，观察生成的 \ce{Ca3(PO4)2} 沉淀。再加入一定量的 \ce{H3PO4}，制取 \ce{CaHPO4}，继续加入一定量的 \ce{H3PO4} 制成 \ce{Ca(H2PO4)2} 溶液。观察实验现象。
\end{enumerate}



\section{胶体的性质}
\subsection*{实验目的}
了解胶体的制备方法及其性质。

\subsection*{实验用品}
烧杯、铁架台、石棉网、试管、试管夹、酒精灯、胶头滴管、量筒、厚纸筒、聚光手电筒。

饱和氯化铁溶液、$1M$ 盐酸、水玻璃、氯化钠溶液、硫酸镁溶液。
\subsection*{实验步骤}
\begin{enumerate}[1.]
  \item 胶体的制备
  \begin{enumerate}[(1)]
    \item 制取氢氧化铁溶胶\quad 用烧杯盛 \qty{30}{ml} 蒸馏水，加热到沸腾。然后逐滴加入饱和氯化铁溶液，边加边振荡，直至溶液变成红褐色，即得氢氧化铁溶胶。
    \item 制取硅酸溶胶
    
    在一个大试管里装入 \qtyrange{5}{10}{ml} 盐酸，并加入 \qty{1}{ml} 的水玻璃，然后用力振荡，即得硅酸的溶胶。
  \end{enumerate}
  \item 胶体的性质
  \begin{enumerate}[(1)]
    \item 胶体的丁达尔现象\quad 用有一小洞的厚纸圆筒（筒径比试管略大些）套在盛有氢氧化铁溶胶（或硅酸溶胶）的试管外面，让日光从小洞进入溶液（或用聚光手电筒照射小孔）。我们从圆筒上方向下观察，可以看到溶胶里有一条光亮的“通路”。用盛有食盐溶液的试管做同样的实验，有没有丁达尔现象发生？
    \item 胶体的凝聚
    \begin{enumerate}[label=\protect\cnum{\arabic*}]
      \item 在一个试管里加入 \qty{3}{ml} 氢氧化铁溶胶，再加入硫酸镁溶液振荡，观察有什么现象发生？
      \item 在一试管里加入 \qty{3}{ml} 氢氧化铁溶胶，然后再加入 \qty{3}{ml} 硅酸溶胶，振荡后，观察有什么现象发生？
      \item 在一试管中加入 \qty{3}{ml} 硅酸溶胶，用试管夹夹住，在酒精灯上加热，观察有什么现象发生？
    \end{enumerate}
  \end{enumerate}
\end{enumerate}
\begin{Theorem}{问题和讨论}
用什么方法可以鉴别一瓶液体是溶液而不是溶胶？
\end{Theorem}


\section{电解质溶液}
\subsection*{实验目的}
巩固电解质溶液的知识，练习 pH 试纸的使用。

\subsection*{实验用品}
试管、试管夹、玻璃棒、酒精灯、pH 试纸。

三种未知 pH 值的溶液、$0.1M$ 盐酸、$0.1M$ 醋酸、硫酸铵溶液、氯化钠溶液、醋酸铵溶液、碳酸钠溶液、醋酸钠溶液、饱和的硫酸铝溶液和碳酸氢钠溶液、酚酞试液、锌粒。

\subsection*{实验步骤}
\begin{enumerate}[1.]
  \item pH 试纸的使用
  
  用干净的玻璃棒沾取少量未知液，滴在一小块 pH 试纸上，观察试纸的颜色变化并跟标准比色卡相比，分别确定三种未知溶液的 pH 值。注意每取过一种未知液后，要用水把玻璃棒冲洗干净，再用来取另一种未知液。比较哪种溶液的氢离子浓度最大。
  \item 强弱电解质
  
  在两个试管里分别 加入 \qty{2}{ml} $0.1M$ 盐酸和 $0.1M$ 醋酸，用玻璃棒沾取少量溶液滴在一小块 pH 试纸上，测定它们的 pH 值。然后分别加入一小粒锌，并加热试管，比较两个试管里反应的快慢。写出反应的离子方程式。
  \item 盐类的水解
  \begin{enumerate}[(2)]
    \item 在四个试管里分别加入 \qty{1}{ml} 碳酸钠、硫酸铵、氯化钠、醋酸铵溶液、用 pH 试纸测定它们的 pH 值。由此可以得出盐类水解的简单规律是什么。写出水解反应的离子方程式。
    \item 在试管里加入 \qty{3}{ml} 醋酸钠溶液，滴入几滴酚酞试剂，观察溶液的颜色。把上述溶液的一半倒入另一个试管里，在酒精灯上加热，比较两个试管里溶液的颜色，并思考温度对水解有什么影响。
    \item 泡沫灭火器的反应原理
    
    在两支试管里分别加入 \qty{2}{ml} 饱和硫酸铝和碳酸氢钠溶液，然后把一个试管里的液体倒入另一个试管里，观察发生的现象。写出反应的离子方程式。
  \end{enumerate}
\end{enumerate}
\begin{Theorem}{问题和讨论}
\begin{enumerate}[1.]
  \item 有人说氯化钠溶液和醋酸铵溶液的 pH 值都接近 7，因此它们都没有发生水解。这种说法对不对？说明理由。
  \item 有下列浓度相同的溶液，把它们按 pH 值由小到大的顺序排列起来。
  \[\ce{Na2S},\quad\ce{CH3COONH4},\quad\ce{HCl},\quad\ce{NaOH},\quad\ce{Al2(SO4)3}\]
\end{enumerate}
\end{Theorem}


\section{中和滴定}
\subsection*{实验目的}
\begin{enumerate}[1.]
  \item 初步学会中和滴定的操作；
  \item 运用和巩固中和滴定的计算技能。
\end{enumerate}
\subsection*{实验用品}
酸式滴定管、碱式滴定管、移液管、滴定管夹、锥形瓶、铁架台、烧杯。

标准的 $0.2N$ \ce{HCl} 溶液\footnote{标准溶液的浓度是用一定方法经过准确测定的。}。待测的 \ce{NaOH} 溶液和 \ce{HCl} 溶液、酚酞试液、甲基橙试液、蒸馏水。

\subsection*{实验步骤}
\begin{enumerate}[1.]
  \item 滴定前的准备工作
  \begin{enumerate}[(1)]
    \item 把滴定管洗涤干净
    \item 检查滴定管是否漏水，酸式滴定管的活塞是否转动灵活，把滴定管分别夹在滴定管夹的两边，并练习从滴定管逐滴滴出液体的操作方法。
    \item 取下酸式滴定管，用标准的 \ce{HCl} 溶液润洗 \numrange{2}{3} 次。每次用酸 \qtyrange{3}{5}{ml}。然后把标准的 \ce{HCl} 溶液注入酸式滴定管刻度 “0” 以上 \qtyrange{2}{3}{cm} 的地方，再把酸式滴定管固定好。在滴定管下放一烧杯，调节到滴定管的尖嘴部分充满溶液，使滴定管内部没有气泡，并使液面处在刻度 “0” 或 “0” 以下的位置。把液面的刻度记入下表。
    \item 采取与(3)相同的方法，把碱式滴定管用待测的 \ce{NaOH} 溶液润洗后装满 \ce{NaOH}溶液，调节好后记下刻度。
  \end{enumerate}
  \item 滴定
  \begin{enumerate}[(1)]
    \item 用移液管取 \qty{25}{ml} 待测的 \ce{NaOH} 溶液，放入锥形瓶里（如无移液管，可用碱式滴定管取 \ce{NaOH} 溶液），滴入 \numrange{2}{3} 滴甲基橙。振荡锥形瓶，使溶液充分混和，这时溶液呈黄色。
    \item 把锥形瓶放在滴定管的下面，瓶下垫一张白纸，小心地滴入酸溶液，边滴边摇动锥形瓶。直到因加入一滴酸溶液后，溶液颜色从黄色变为橙色，并不立即褪去为止。这就表示已到滴定终点。最后记下滴定管液面所在的刻度，填入下表。
    \item 把锥形瓶里的溶液倒掉，用蒸馏水把锥形瓶洗涤干净。按上述操作重作一次，并把滴定前后液面刻度分别记入下表。
    \begin{tablehere}
      \begin{tblr}{colspec={crX[c]X[c]X[r]},hline{3}=0.8pt,row{1,2}={m,c}}
        \SetCell[r=2]{m,c} 滴定次数 & \SetCell[r=2]{m,c} 待测碱溶液体积 & \SetCell[c=3]{m,c} 标准酸溶液的体积 & & \\
        & & 滴定前刻度 & 滴定后刻度 & 体积\\
        第一次 & （\unit{ml}）& & & （\unit{ml}）\\
        第二次 & （\unit{ml}）& & & （\unit{ml}）\\
      \end{tblr}
    \end{tablehere}
    \item 取两次测定数值的平均值计算待测NaOH溶液的当量浓度。
  \end{enumerate}
  \item 用已经算出准确浓度的 \ce{NaOH} 溶液,滴定另一未知浓度的 \ce{HCl} 溶液。
  
  取 \qty{25}{ml} 未知浓度的 \ce{HCl} 溶液放入锥形瓶里，加入几滴酚酞试液作指示剂用碱式滴定管滴入已知浓度的 \ce{NaOH} 溶液。滴到溶液由无色变成红色，并在半分钟内红色不褪为止。再重复滴定操作一次，记下数据。最后计算出这种 \ce{HCl} 溶液的浓度。
\end{enumerate}
\begin{Theorem}{问题和讨论}
\begin{enumerate}[1.]
  \item 锥形瓶是否需要用被滴定的溶液润洗？为什么？
  \item 观察滴定管里液面的高度时应注意什么？
\end{enumerate}
\end{Theorem}


\section{中和热的测定}
\subsection*{实验目的}
\begin{enumerate}[1.]
  \item 认识中和反应是放热反应；
  \item 测定强酸跟强碱反应的中和热。
\end{enumerate}
\subsection*{实验用品}
大烧杯（\qty{500}{ml}）、小烧杯（\qty{100}{ml}）、温度计、量筒（\qty{50}{ml}）两个、碎纸条、硬纸板（中心有一小孔）、玻璃棒。

$1M$ 盐酸、$1.1M$ 氢氧化钠溶液\footnote{为了保证 $1M$ 的盐酸完全被氢氧化钠中和，而采用了 $1.1M$ 的氢氧化钠溶液，使碱稍稍过量。}。

\subsection*{实验步骤}
\begin{enumerate}[1.]
  \item 在大烧杯中放入一些碎纸条，把小烧杯放在碎纸条上。然后再在大小烧杯之间填满碎纸条（\cref{fig:a5}），大烧杯上用硬纸板盖好，以达到保温、隔热、减少实验过程中热量损失的目的。
  \begin{figure}
    \includegraphics{A-5.pdf}
    \caption{中和热的测定}\label{fig:a5}
  \end{figure}
  \item 用一个量筒量取 \qty{50}{ml} $1M$ 的盐酸，倒入小烧杯中。并用温度计测量盐酸的温度，记入下表。然后把温度计上的酸用水冲洗干净。
\begin{tablehere}
  \begin{minipage}{\linewidth}
  \begin{tblr}{colspec={c*3{X[c]}cc},hline{3}=0.8pt}
    \SetCell[r=2]{m,c} 实验次数
    &\SetCell[c=3]{m,c} 起始温度（\unit{\celsius}） & & &\SetCell[r=2]{m,c} {终止温度（\unit{\celsius}） \\ $t_2$}& 温度差（\unit{\celsius}）\\
    &\ce{HCl} & \ce{NaOH} & 平均值 $t_1$ & & $t_2-t_1$\\
    1 & &  &  &  & \\
    2 & &  &  &  & \\
    3 & &  &  &  & \\
  \end{tblr}
\end{minipage}
\end{tablehere}
  \item 用另一个量筒量取 \qty{50}{ml} $1.1M$ 的氢氧化钠溶液，并用温度计测量氢氧化钠溶液的温度，记入下表。
  \item 把温度计放入小烧杯的盐酸中，并把量筒中的氢氧化钠溶液一次倒入小烧杯（注意不要洒到外面）。用玻璃棒轻轻搅动溶液，并准确读取混和溶液的最高温度，记为终止温度。记入下表。
  \item 根据实验数据计算中和热
  
  为了使计算简便一些，我们近似地认为：
  \begin{enumerate}[(1)]
    \item $1M$ 盐酸和 $1.1M$ 氢氧化钠溶液的密度都是 \qty{1}{g/ml}，所以 \qty{50}{ml} $1M$ 盐酸的质量 $m_1=\qty{50}{g}$，\qty{50}{ml} $1.1M$ 氢氧化钠的质量是 $m_2=\qty{50}{g}$。
    \item 中和后生成的溶液的比热是 $C=\qty{1}{Cal.g^{-1}.K^{-1}}$。由此，可以计算出 \qty{50}{ml} $1M$ 盐酸跟 \qty{50}{ml} $1.1M$ 氢氧化钠发生中和反应时，放出的热量 $Q$ 为：
    \begin{align*}
      Q &=(m_1+m_2)\cdot C \cdot (t_2-t_1)\\
      &=100(t_2-t_1)\,\unit{Cal}
    \end{align*}
    又因 \qty{50}{ml} $1M$ 盐酸中含有 $\dfrac{50}{1000} = \qty{0.05}{mol}$ 的 \ce{HCl}，\qty{0.05}{mol} 的 \ce{HCl} 跟 \qty{0.05}{mol} 氢氧化钠发生中和反应，生成 \qty{0.05}{mol} 的水。生成 \qty{0.05}{mol} 的水，放出的热量是 $Q$。所以生成 \qty{1}{mol} 水时，放出的热量，即为中和热。
    \[\text{中和热}=\frac{Q}{0.05}=\frac{100(t_2-t_1)}{0.05}\,\unit{Cal}\]
  \end{enumerate}
\end{enumerate}
\begin{Theorem}{问题和讨论}
要想提高测量的准确程度，实验中应该注意什么？
\end{Theorem}


\section{原电池\texorpdfstring{\quad}{ }金属的电化腐蚀}
\subsection*{实验目的}
\begin{enumerate}[1.]
  \item 了解原电池的原理；
  \item 认识金属电化腐蚀的原因。
\end{enumerate}
\subsection*{实验用品}
烧杯、试管、滴管、导线、电流计、小刀。

稀硫酸、稀盐酸、硫酸铜溶液、食盐溶液、赤血盐溶液、镀锌铁片、镀锡铁片、锌片、铜片、锌粒、不纯净的锌。

\subsection*{实验步骤}
\begin{enumerate}[1.]
  \item 原电池原理
  
  把一块纯净的锌片插入盛有稀硫酸的烧杯里，观察发生的现象；再插入一块铜片，观察铜片上有没有气体放出？用导线把锌片和铜片连接起来，观察铜片上有没有气体放出？在导线中间接入一个电流计，观察指针是否偏转？解释原因，并写出正负极上的反应。
  \item 金属的电化腐蚀
  \begin{enumerate}[(1)]
    \item 用锋利的小刀分别在一小块镀锌铁（白铁）和镀锡铁（马口铁）上刻划，露出里面的铁来，然后在刻痕上滴一滴食盐溶液，再滴一滴赤血盐\footnote{赤血盐（铁氰化钾）溶液遇到 \ce{Fe^2+} 会产生蓝色沉淀。}溶液。静置 \qty{5}{min}，观察发生的现象，如果现象不明显，再静置 \qty{10}{min}，观察现象。静置时，可以先进行下边的实验。
    \item 在一个试管里放一小块纯净的锌，在另一个试管里放一小块含有杂质的锌\footnote{或用在硫酸铜溶液中浸过的锌（即锌的表面上附着铜）。},然后在两个试管里各注入稀硫酸 \qty{2}{ml}。观察发生的现象，哪个试管里的反应进行得较为剧烈？为什么？
    \item 在上面实验（2）盛着纯净的锌和稀硫酸的试管里，加入少量的硫酸铜溶液。观察有什么变化。为什么？
  \end{enumerate}

\end{enumerate}
\begin{Theorem}{问题和讨论}
\begin{enumerate}[1.]
  \item 实验室制取氢气时，选用纯净的锌好，还是选用含有少量杂质的锌好？
  \item 根据你做的实验，小结组成原电池装置必须具备的几个条件。
\end{enumerate}
\end{Theorem}


\section{电解\texorpdfstring{\quad}{ }电镀}
\subsection*{实验目的}
\begin{enumerate}[1.]
  \item 巩固对电解、电镀原理的认识；
  \item 了解无氰镀锌的简单操作。
\end{enumerate}
\subsection*{实验用品}
铁制镀件、砂纸、烧杯、 U 形管、玻璃棒、导线、 直流电源、石墨电极、碘化钾淀粉试纸。

$1N$ 氢氧化钠溶液、锌片、电镀液、钝化液\footnote{电镀液和钝化液都由教师配制，实验时发给学生。}、20\% 盐酸、氯化铜溶液。

\subsection*{实验步骤}
\begin{enumerate}[1.]
  \item 电解
  
  在 U 形管中，注入氯化铜溶液，插入两根石墨电极，接通电源，把湿润的碘化钾淀粉试纸放在阳极碳棒附近检验放出的气体并观察管内发生的现象。
  \item 电镀锌
  \begin{enumerate}[(1)]
    \item 镀件除油、除锈
    
    把镀件用砂纸打磨干净，放在 $1N$ 的氢氧化钠溶液里，除去油污，然后用清水洗净。再放入 20\% 盐酸中除锈，几分钟后取出，并用清水洗净。
    \item 电镀
    
    把镀件接在电压为 \qtyrange{2}{3}{V} 的直流电源的负极上，锌片接到正极上。把两极平行浸入电镀液中，两极之间距离约为 \qty{5}{cm}，\qty{5}{10}{min} 后取出，镀件表面便镀了一层锌。
    \item 钝化
    把镀件取出，用清水洗净，再放入钝化液中进行钝化\footnote{钝化是使镀层表面生成一层钝化膜（结构紧密的氧化物）以保护镀层。}。\qtyrange{5}{10}{s} 后取出，用清水洗净，便可看到镀件表层出现彩色花纹。
  \end{enumerate}
\end{enumerate}
\begin{Theorem}{问题和讨论}
\begin{enumerate}[1.]
  \item 在电解水的过程里，常加入氢氧化钠溶液或硫酸，这是为什么？
  \item 现有一个烧杯，里面盛有氯化铜溶液，还有两根导线、两个石墨电极，你能用以上仪器和药品设计一个简单的实验来判断一个电池的正负极吗？说明理由。
\end{enumerate}
\end{Theorem}


\section{铝和氢氧化铝的化学性质}
\subsection*{实验目的}
认识铝的主要化学性质和氢氧化铝的两性。

\subsection*{实验用品}
试管、试管夹、酒精灯、滤纸。

稀盐酸、稀硫酸、浓硝酸、氨水、\ce{NaOH} 溶液、\ce{KOH} 溶液、\ce{Hg(NO3)2} 溶液、\ce{Al2(SO4)3} 溶液、 \ce{AlCl3} 溶液、铝条（或铝片）、蒸馏水。

\subsection*{实验步骤}
\begin{enumerate}[1.]
  \item 铝的化学性质
  \begin{enumerate}[(1)]
    \item 铝跟酸的反应\quad 在三个试管里各放入大小相近的铝条，然后往三个试管里分别注入 \qtyrange{2}{3}{ml} 稀盐酸、稀硫酸和浓硝酸 （所用试管必须干燥）。注意观察哪个试管里有反应发生，哪个试管里没有发生变化。写出化学方程式和离子方程式。
    \item 铝跟碱的反应\quad 在两个试管里各放入大小相近的铝条，然后分别注入 \qtyrange{2}{3}{ml} \ce{NaOH} 溶液和 \ce{KOH} 溶液，稍加热，观察发生的现象。写出这两个反应的化学方程式和离子方程式。
    \item 氧化铝薄膜的保护作用\quad 把铝条放入盛有 \qtyrange{2}{3}{ml} \ce{NaOH} 溶液的试管里，稍加热，使铝条跟 \ce{NaOH} 溶液起反应，经过 \qtyrange{2}{3}{min} 后，把铝条取出，用蒸馏水冲洗干净。然后，把铝条插入 \ce{Hg(NO3)2} 溶液里浸泡，铝跟 \ce{Hg(NO3)2} 发生了置换反应，当观察到铝条上有少量汞析出时，取出铝条，用碎滤纸轻轻地把它的表面吸干，把铝条放在干燥的滤纸上。观察铝条表面发生的现象。写出有关的化学方程式。
  \end{enumerate}
  \item 氢氧化铝的两性
  \begin{enumerate}[(1)]
    \item 氢氧化铝的生成\quad 用三个试管各取 \qty{3}{ml} \ce{Al2(SO4)3} 溶液，往这三个试管里分别滴加 \ce{NaOH} 溶液、\ce{KOH} 溶液和氨水，直到产生大量的沉淀为止。写出有关的化学方程式和离子方程式。

再用 \ce{AlCl3} 溶液代替 \ce{Al2(SO4)3} 溶液重复做上面的实验。写出有关的化学方程式和离子方程式。
    \item 氢氧化铝跟酸和碱的反应\quad 取两个盛有 \ce{Al(OH)3} 沉淀的试管，分别注入稀盐酸和强碱溶液（\ce{NaOH} 溶液或 \ce{KOH} 溶液），振荡。观察发生的现象。写出有关的化学方程式和离子方程式。

另取一个盛有 \ce{Al(OH)3} 沉淀的试管，往这个试管里注入氨水，振荡。观察发生的现象。
  \end{enumerate}
\end{enumerate}
\begin{Theorem}{问题和讨论}
\begin{enumerate}[1.]
  \item 铝跟稀盐酸的反应与跟浓硝酸的反应有什么不同？为什么？
  \item 比较往 \ce{Al(OH)3} 沉淀里加入强碱溶液和加入氨水有什么不同？为什么？
  \item 往 \ce{AlCl3} 溶液里逐滴加入 \ce{NaOH} 溶液直至过量，和往 \ce{NaOH} 溶液里逐滴加入 \ce{AlCl3} 溶液直至过量的反应的现象相同吗？
\end{enumerate}
\end{Theorem}


\section{分子量的测定}
\subsection*{实验目的}
了解一种测定分子量的方法——蒸气密度法。

\subsection*{实验用品}
托盘天平、圆底烧瓶（\qty{250}{ml}）、烧杯（\qty{500}{ml}）、量筒 （\qty{10}{ml}、\qty{100}{ml} 各一个）、铝箔、棉线、针、酒精灯、 铁架台、铁圈、石棉网。

四氯化碳。
\subsection*{实验步骤}
\begin{enumerate}[1.]
  \item 沸水的准备\quad 往 \qty{500}{ml} 烧杯中加入约 \qty{200}{ml} 热水，放在铁架台的石棉网上，用酒精灯加热到沸腾（\qty{100}{\celsius}）。供下面实验使用。
  \item 称量烧瓶的质量\quad 取一干燥而洁净的 \qty{250}{ml} 的圆底烧瓶，与封口用的铝箔、棉线一起用托盘天平称量其质量 $W_1$（准确到 \qty{0.1}{g}）。
  \item 四氯化碳的气化\quad 往上述已称量过的烧瓶中加入约 \qty{2}{ml} 四氯化碳，用铝箔和棉线将瓶口封好。用针在铝箔上扎一小孔。然后将烧瓶放入盛有沸水的烧杯中（\cref{fig:a6}），并使烧瓶没入水中（此时要继续加热，使水保持沸腾状态）烧瓶中的四氯化碳受热气化，同时将瓶中空气排出。
  \begin{figure}
    
\includegraphics{A-6.pdf}
    \caption{四氯化碳的气化}\label{fig:a6}
  \end{figure}
  \item 四氯化碳质量的称量\quad 待四氯化碳完全气化后，将烧瓶从沸水中取出。瓶中的四氯化碳蒸气因被空气冷却而全部凝成液体。当烧瓶冷到室温后，将其外部擦干，并称量其质量 $W_2$。
  \item 烧瓶体积的测量\quad 将称量过的烧瓶中的液态四氯化碳倒出，然后用 \qty{100}{ml} 量筒取水，倒入烧瓶中，当烧瓶被水充满时，水的体积即为烧瓶的体积 $V$（也可先用水将烧瓶充满，然后将水倒在量筒中，量出水的体积）。
  \item 四氯化碳分子量的计算
  
  根据 
  \[PV=nRT\]
  \begin{EqDesc}{P}
    \item[P] 近似地看作 \qty{1}{atm}
    \item[T] $273 + 100 = \qty{373}{K}$
    \item[R] \qty{0.082}{L.atm.K^{-1}.mol^{-1}}
    \item[n] $n=\dfrac{W_2-W_1}{\mathrm{M}}$（$\mathrm{M}$：四氯化碳的摩尔质量）
  \end{EqDesc}
  \begin{align*}
    \mathrm{M} & =\frac{RT(W_2-W_1)}{PV} \\
     & =\frac{0.082\times 373(W_2-W_1)}{V}
  \end{align*}
  摩尔质量的数值与分子量是相同的。
\end{enumerate}
\begin{Theorem}{问题和讨论}
\begin{enumerate}[1.]
  \item 蒸气密度法测分子量的方法 能够用来测定遇热不挥发或遇热分解的物质吗？
  \item 试讨论下列情况对实验结果有什么影响？
  \begin{enumerate}[(1)]
    \item 压强低于 \qty{1}{atm}。
    \item 水未沸腾。
    \item 四氯化碳气化不完全。
  \end{enumerate}
\end{enumerate}
\end{Theorem}


\section{实验习题}
\subsection*{实验目的}
巩固已学知识，培养实验能力。

\subsection*{实验习题}
\begin{enumerate}[1.]
  \item 怎样证明蔗糖是非电解质，醋酸是弱电解质，食盐是强电解质？
  \item 电解硫酸钠溶液，说明阴极、阳极上发生了什么反应（在这里，\ce{OH-} 比 \ce{SO4^2-} 易失去电子）。
  \item 检验下列水溶液的酸碱性，说明原因。
  \begin{tasks}[item-indent=5em,label=\Alph*.,label-width=15pt](2)
    \task \ce{NaNO3}，
    \task \ce{K2CO3}，
    \task \ce{NH4Cl}，
    \task \ce{CH2COONH4}。
  \end{tasks}
  \item 在试管里放入少量明矾，加水溶解，试验它的酸碱性. 判断它是不是胶体。
  \item 判断下列物质是不是胶体。
  \begin{tasks}[item-indent=5em,label=\Alph*.,label-width=15pt](3)
    \task 食盐水，
    \task 蔗糖水，
    \task 淀粉溶液，
    \task*(2) 用水稀释的蓝墨水，
    \task 一瓶未知液。
  \end{tasks}
  \item 有三瓶水，分别是蒸馏水、永久硬水和暂时硬水，通过实验确定哪个瓶子里装的是哪种水。
  \item 有 \ce{MgSO4}、\ce{BaCl2}、\ce{Al2(SO4)3}、\ce{FeCl3} 四瓶溶液，通过实验把它们鉴别出来。
  \item 有五瓶溶液,已知它们分别是 \ce{KCl}、\ce{Ba(NO3)2}、\ce{Na2CO3}、\ce{Na2SO4} 和 \ce{FeCl3} 溶液，不用任何其它试剂和试纸，通过观察和实验确定哪个瓶子里装的是哪种溶液。
  \item 只用一种试剂，鉴别下列物质的溶液。
  \[\ce{NaCl},\quad\ce{MgCl2},\quad\ce{Al2(SO4)3},\quad\ce{NH4Cl}\]
  \item 用已知浓度的盐酸滴定未知浓度的氢氧化钾溶液，测定它的浓度。
\end{enumerate}